% !TeX root = ../main.tex
\Needspace{16\baselineskip}
\section{Prodotto di matrici e trasposizione}
\begin{definizione}
Se $A\in M(m,n,K)$ e $B\in M(n,p,K)$, il \emph{prodotto righe per colonne} $AB\in M(m,p,K)$ è definito da
\[
 (AB)_{ij}=\sum_{h=1}^n a_{ih}b_{hj}.
\]
Il numero di colonne del primo fattore deve coincidere con il numero di righe del secondo.
\end{definizione}
\begin{schema}
\begin{tikzpicture}[x=1cm,y=1cm,font=\small]
 % Ogni tabella è una matrice schematica: le righe estreme sono omesse.
 \fill[coloreschema!13] (-.42,-.28) rectangle (2.12,.28);
 \fill[coloreDefinizioni!13] (4.15,-.98) rectangle (4.85,.98);
 \fill[coloreschema!18] (7.35,-.28) rectangle (8.05,.28);
 \foreach \xx/\top/\mid/\bot in {0/\vdots/a_{i1}/\vdots,.85/\vdots/\cdots/\vdots,1.7/\vdots/a_{in}/\vdots}{
  \node at (\xx,.7) {$\top$};\node at (\xx,0) {$\mid$};\node at (\xx,-.7) {$\bot$};}
 \foreach \yy/\entry in {.7/b_{1j},0/\vdots,-.7/b_{nj}}{
  \node at (3.65,\yy) {$\cdots$};\node at (4.5,\yy) {$\entry$};\node at (5.35,\yy) {$\cdots$};}
 \foreach \xx in {6.85,7.7,8.55}{\node at (\xx,.7) {$\vdots$};\node at (\xx,-.7) {$\vdots$};}
 \node at (6.85,0) {$\cdots$};\node at (7.7,0) {$c_{ij}$};\node at (8.55,0) {$\cdots$};
 \foreach \left/\right in {-.52/2.22,3.13/5.87,6.33/9.07}{
  \draw (\left+.12,1.02)--(\left,1.02)--(\left,-1.02)--(\left+.12,-1.02);
  \draw (\right-.12,1.02)--(\right,1.02)--(\right,-1.02)--(\right-.12,-1.02);}
 \node at (2.67,0) {$\cdot$};\node at (6.1,0) {$=$};
 \node at (.85,1.45) {$A\ (m\times n)$};
 \node at (4.5,1.45) {$B\ (n\times p)$};
 \node at (7.7,1.45) {$C=AB\ (m\times p)$};
 \node[text=coloreschema] at (.85,-1.4) {riga $i$};
 \node[text=coloreDefinizioni] at (4.5,-1.4) {colonna $j$};
 \node at (7.7,-1.4) {posizione $(i,j)$};
\end{tikzpicture}

$c_{ij}=a_{i1}b_{1j}+\cdots+a_{in}b_{nj}$: si moltiplicano le entrate corrispondenti della riga $i$ di $A$ e della colonna $j$ di $B$, poi si sommano i prodotti. Lo schema omette le altre righe e colonne; il risultato si chiama $c_{ij}$ perché appartiene a $C=AB$, mentre $a_{ij}$ indica un'entrata di $A$.
\end{schema}

Per una colonna $x=\left(\begin{smallmatrix}x_1\\\vdots\\x_n\end{smallmatrix}\right)$, la componente $i$ di $Ax$ è $\sum_j a_{ij}x_j$. Un sistema lineare si scrive dunque $Ax=b$. Se $A_1,\ldots,A_n$ sono le colonne di $A$, vale inoltre
\[
 Ax=x_1A_1+\cdots+x_nA_n.
\]
Il prodotto di una riga $(a_1,\ldots,a_n)$ per una colonna $\left(\begin{smallmatrix}b_1\\\vdots\\b_n\end{smallmatrix}\right)$ è lo scalare $\sum_h a_hb_h$; su $\R$ è il prodotto scalare usuale.

\begin{esempio}
Per
\[
 A=\begin{pmatrix}2&1\\-3&3\end{pmatrix},\qquad
 B=\begin{pmatrix}1&4\\2&-1\end{pmatrix}
\]
si ottiene
\[
 AB=\begin{pmatrix}4&7\\3&-15\end{pmatrix},\qquad
 BA=\begin{pmatrix}-10&13\\7&-1\end{pmatrix}.
\]
Per esempio, $c_{12}=2\cdot4+1\cdot(-1)=7$: si usa la prima riga di $A$ e la seconda colonna di $B$. Anche quando entrambi i prodotti sono definiti, in generale $AB\ne BA$.
\end{esempio}

\textbf{La matrice identità.} Con $I_n$ indichiamo la matrice quadrata di tipo $n\times n$ definita da
\[
 I_n=\begin{pmatrix}
 1&0&\cdots&0\\0&1&\cdots&0\\
 \vdots&\vdots&\ddots&\vdots\\0&0&\cdots&1
 \end{pmatrix},\qquad
 (I_n)_{ij}=\begin{cases}1,&i=j,\\0,&i\ne j.\end{cases}
\]
Il pedice $n$ ne specifica il numero di righe e di colonne: per esempio $I_2=\left(\begin{smallmatrix}1&0\\0&1\end{smallmatrix}\right)$. Moltiplicare per l'identità conserva la matrice, purché le dimensioni siano compatibili.

\begin{proposizione}
Per matrici di dimensioni compatibili valgono:
\begin{enumerate}[label=\arabic*.]
\item l'associatività $(AB)C=A(BC)$;
\item le distributività $A(B+C)=AB+AC$ e $(A+B)C=AC+BC$;
\item le identità $AI_n=A$ e $I_mA=A$ per $A\in M(m,n,K)$.
\end{enumerate}
\end{proposizione}
\begin{proof}
\begin{enumerate}[label=\arabic*.]
\item Se $A$ è $m\times n$, $B$ è $n\times p$ e $C$ è $p\times q$, allora
\begin{align*}
 ((AB)C)_{ij}
 &=\sum_{\ell=1}^p\left(\sum_{h=1}^n a_{ih}b_{h\ell}\right)c_{\ell j}
 =\sum_{h=1}^n\sum_{\ell=1}^p a_{ih}b_{h\ell}c_{\ell j}\\
 &=\sum_{h=1}^n a_{ih}\left(\sum_{\ell=1}^p b_{h\ell}c_{\ell j}\right)
 =(A(BC))_{ij}.
\end{align*}
\item Per la prima distributività, $\sum_h a_{ih}(b_{hj}+c_{hj})=\sum_h a_{ih}b_{hj}+\sum_h a_{ih}c_{hj}$; l'altra segue allo stesso modo.
\item Nel prodotto con l'identità sopravvive soltanto il termine con indice coincidente, che è l'entrata corrispondente di $A$.
\end{enumerate}
\end{proof}

\begin{definizione}
La \emph{trasposta} di $A\in M(m,n,K)$ è $A^T\in M(n,m,K)$, definita da $(A^T)_{ij}=a_{ji}$: le righe diventano colonne e viceversa.
\end{definizione}
\begin{esempio}
\[
 A=\begin{pmatrix}2&1&3&0\\4&-2&1&3\end{pmatrix}
 \quad\Longrightarrow\quad
 A^T=\begin{pmatrix}2&4\\1&-2\\3&1\\0&3\end{pmatrix}.
\]
\end{esempio}
\begin{proposizione}
Valgono le seguenti proprietà:
\begin{enumerate}[label=\arabic*.]
\item $(A^T)^T=A$ e $(A+B)^T=A^T+B^T$;
\item $(AB)^T=B^TA^T$.
\end{enumerate}
\end{proposizione}
\begin{proof}
\begin{enumerate}[label=\arabic*.]
\item Le prime due uguaglianze seguono scambiando gli indici una o due volte.
\item Per il prodotto,
\[
 ((AB)^T)_{ij}=(AB)_{ji}=\sum_h a_{jh}b_{hi}
 =\sum_h (B^T)_{ih}(A^T)_{hj}=(B^TA^T)_{ij}.
\]
La trasposizione inverte quindi l'ordine dei fattori.
\end{enumerate}
\end{proof}

\section{Matrici inverse}
\begin{definizione}
Per una matrice quadrata $A\in M(n,n,K)$, una matrice $B$ è un'\emph{inversa destra} se $AB=I_n$, un'\emph{inversa sinistra} se $BA=I_n$. È un'\emph{inversa} se valgono entrambe le identità; in questo caso $A$ è detta invertibile.
\end{definizione}
\begin{proposizione}
Se $B$ è inversa sinistra di $A$ e $C$ è inversa destra, allora $B=C$. In particolare l'inversa, quando esiste, è unica e si indica con $A^{-1}$.
\end{proposizione}
\begin{proof}
Per associatività, $B=B(AC)=(BA)C=C$.
\end{proof}
Se $A$ è invertibile, il sistema $Ax=b$ ha l'unica soluzione $x=A^{-1}b$: moltiplicando per $A^{-1}$ si ottiene necessariamente questa espressione e, viceversa, $A(A^{-1}b)=b$.

% !TeX root = ../main.tex
\subsection{Le operazioni sulle righe come prodotti di matrici}
Il procedimento con il blocco identità usato per conservare le relazioni fra generatori ha una spiegazione generale. Per operare sulle righe di una matrice $A\in M(m,n,K)$ si moltiplica \emph{a sinistra} per opportune matrici quadrate di ordine $m$.

Indichiamo con $E_{ij}\in M(m,m,K)$ la matrice che ha $1$ nella posizione $(i,j)$ e $0$ in tutte le altre. Per $i\ne j$ le tre \emph{matrici elementari} corrispondenti alle operazioni di Gauss sono
\begin{align*}
 P_{ij}&=I_m-E_{ii}-E_{jj}+E_{ij}+E_{ji},\\
 M_i(\lambda)&=I_m+(\lambda-1)E_{ii}\quad(\lambda\ne0),\\
 S_{ij}(\alpha)&=I_m+\alpha E_{ij}.
\end{align*}
Si ottengono da $I_m$ effettuando, rispettivamente, lo scambio delle righe $i,j$, la moltiplicazione della riga $i$ per $\lambda$, e la sostituzione $R_i\leftarrow R_i+\alpha R_j$. Per esempio, in ordine $3$,
\[
 P_{12}=\begin{pmatrix}0&1&0\\1&0&0\\0&0&1\end{pmatrix},\quad
 M_2(\lambda)=\begin{pmatrix}1&0&0\\0&\lambda&0\\0&0&1\end{pmatrix},\quad
 S_{21}(\alpha)=\begin{pmatrix}1&0&0\\\alpha&1&0\\0&0&1\end{pmatrix}.
\]
\begin{proposizione}
\begin{enumerate}[label=\arabic*.]
\item Il prodotto $P_{ij}A$ scambia le righe $i,j$ di $A$; $M_i(\lambda)A$ moltiplica la riga $i$ per $\lambda$; $S_{ij}(\alpha)A$ somma alla riga $i$ il multiplo $\alpha$ della riga $j$. Le altre righe restano invariate.
\item Tutte queste matrici sono invertibili:
\[
 P_{ij}^{-1}=P_{ij},\qquad M_i(\lambda)^{-1}=M_i(1/\lambda),\qquad
 S_{ij}(\alpha)^{-1}=S_{ij}(-\alpha).
\]
\end{enumerate}
\end{proposizione}
\begin{proof}
\begin{enumerate}[label=\arabic*.]
\item La riga $h$ di un prodotto $EA$ è la combinazione delle righe di $A$ i cui coefficienti sono le entrate della riga $h$ di $E$. Una riga uguale alla riga $h$ di $I_m$ riproduce dunque la riga $h$ di $A$. Le righe modificate nelle tre matrici danno precisamente le tre operazioni indicate.
\item Le formule delle inverse corrispondono a ripetere lo scambio, moltiplicare per $1/\lambda$ e sottrarre lo stesso multiplo: in ciascun caso si torna alla matrice iniziale.
\end{enumerate}
\end{proof}

\subsection{Registrare tutte le operazioni nel blocco destro}
Se le operazioni elementari sono rappresentate, nell'ordine, da $E_1,\ldots,E_s$, il risultato è
\[
 A'=E_s\cdots E_2E_1A=BA,\qquad B=E_s\cdots E_2E_1.
\]
La prima operazione è il fattore più vicino ad $A$; la matrice $B$ è invertibile, con $B^{-1}=E_1^{-1}E_2^{-1}\cdots E_s^{-1}$: i fattori si annullano a coppie in ordine inverso. Per calcolare $B$ non occorre moltiplicare separatamente tutti i fattori: basta ridurre il blocco sinistro di $(A\mid I_m)$, effettuando le stesse operazioni sull'intera riga. Infatti
\[
 B(A\mid I_m)=(BA\mid B)=(A'\mid B).
\]
\begin{schema}
\begin{tikzpicture}[font=\small]
 \node (l) at (0,0) {$\displaystyle(A\mid I_m)$};
 \node (r) at (7,0) {$\displaystyle(A'\mid B)$};
 \draw[-{Stealth},coloreschema] (l)--node[above] {stesse operazioni su entrambi i blocchi} (r);
 \node at (0,-.7) {coefficienti \quad registro iniziale};
 \node at (7,-.7) {forma a scala \quad operazioni accumulate};
\end{tikzpicture}
\end{schema}

\Needspace{14\baselineskip}
\begin{esempio}
Cerchiamo una matrice $B$ tale che $BA$ sia una forma a scala di
\[
 A=\begin{pmatrix}1&3&-2\\2&5&1\\-1&0&3\end{pmatrix}.
\]
Partiamo da $(A\mid I_3)$ e applichiamo $R_2\leftarrow R_2-2R_1$ e $R_3\leftarrow R_3+R_1$:
\[
 \left(\begin{array}{rrr|rrr}1&3&-2&1&0&0\\2&5&1&0&1&0\\-1&0&3&0&0&1\end{array}\right)
 \sim\left(\begin{array}{rrr|rrr}1&3&-2&1&0&0\\0&-1&5&-2&1&0\\0&3&1&1&0&1\end{array}\right).
\]
Poi $R_3\leftarrow R_3+3R_2$ dà
\[
 \left(\begin{array}{rrr|rrr}1&3&-2&1&0&0\\0&-1&5&-2&1&0\\0&0&16&-5&3&1\end{array}\right).
\]
Quindi
\[
 B=\begin{pmatrix}1&0&0\\-2&1&0\\-5&3&1\end{pmatrix},\qquad
 BA=\begin{pmatrix}1&3&-2\\0&-1&5\\0&0&16\end{pmatrix}=A'.
\]
L'ultima uguaglianza si verifica anche direttamente con il prodotto righe per colonne. La forma a scala non è unica: questa $B$ registra la particolare sequenza di operazioni scelta.
\end{esempio}

\subsection{Calcolare l'inversa con il blocco identità}
Per una matrice quadrata $A\in M(n,n,K)$ proseguiamo la riduzione fino alla forma a scala ridotta. Se il blocco sinistro diventa $I_n$, il blocco destro è l'inversa:
\[
 (A\mid I_n)\sim(I_n\mid B),\qquad BA=I_n\quad\Longrightarrow\quad B=A^{-1}.
\]
Infatti $B$ è un prodotto di matrici elementari, dunque è invertibile. Moltiplicando $BA=I_n$ a sinistra per $B^{-1}$ si trova $A=B^{-1}$, e quindi anche $AB=I_n$. Se invece la riduzione produce una riga nulla nel blocco sinistro, il sistema omogeneo $Ax=0$ ha una variabile libera e una soluzione non nulla: un'inversa darebbe $x=A^{-1}0=0$, una contraddizione. Il criterio in termini di rango sarà ripreso nella sezione~\ref{sec:invertibilita-applicazioni}.

\begin{esercizio}
Verificare se la matrice
\[
 A=\begin{pmatrix}2&-1&3\\-6&2&-5\\0&1&-3\end{pmatrix}
\]
è invertibile e, in caso affermativo, calcolarne l'inversa.
\end{esercizio}
\begin{svolgimento}
\textbf{Riduzione a scala.} Applichiamo $R_2\leftarrow R_2+3R_1$ e poi $R_3\leftarrow R_3+R_2$:
\begin{align*}
 (A\mid I_3)&=
 \left(\begin{array}{rrr|rrr}2&-1&3&1&0&0\\-6&2&-5&0&1&0\\0&1&-3&0&0&1\end{array}\right)\\
 &\sim\left(\begin{array}{rrr|rrr}2&-1&3&1&0&0\\0&-1&4&3&1&0\\0&1&-3&0&0&1\end{array}\right)
 \sim\left(\begin{array}{rrr|rrr}2&-1&3&1&0&0\\0&-1&4&3&1&0\\0&0&1&3&1&1\end{array}\right).
\end{align*}
\textbf{Eliminazione sopra i pivot.} Con $R_2\leftarrow R_2-4R_3$ e $R_1\leftarrow R_1-3R_3$ si ottiene
\[
 \left(\begin{array}{rrr|rrr}2&-1&0&-8&-3&-3\\0&-1&0&-9&-3&-4\\0&0&1&3&1&1\end{array}\right).
\]
Cambiamo segno a $R_2$, sommiamola a $R_1$ e dividiamo infine $R_1$ per $2$:
\[
 \left(\begin{array}{rrr|rrr}2&0&0&1&0&1\\0&1&0&9&3&4\\0&0&1&3&1&1\end{array}\right)
 \sim\left(\begin{array}{rrr|rrr}1&0&0&1/2&0&1/2\\0&1&0&9&3&4\\0&0&1&3&1&1\end{array}\right).
\]
La matrice è dunque invertibile e
\[
 A^{-1}=\begin{pmatrix}1/2&0&1/2\\9&3&4\\3&1&1\end{pmatrix}.
\]
Moltiplicando si verifica $AA^{-1}=A^{-1}A=I_3$: questo controlla il risultato indipendentemente dalla sequenza di riduzione.
\end{svolgimento}

